\documentclass[a4paper,landscape,12pt]{extarticle}
\usepackage[left=0.5in, right=0.5in, top=0.4in, bottom=0.55in]{geometry}
\usepackage{graphicx} % Required for inserting images
\usepackage{amsmath}
\usepackage{amssymb}
\usepackage{siunitx}
\usepackage{caption}
\usepackage{fancyhdr}
\usepackage{enumitem}
\usepackage{ulem}

\pagestyle{fancy}
\fancyhf{}
\renewcommand{\headrulewidth}{0pt}
\fancyfoot[R]{\Large\thepage\ / \pageref{LastPage}}

% Header style for slides
\newcommand{\slideheader}[2]{%
\noindent
{\Large\textbf{#1}}\hfill {\large #2}
\par\vspace{0.4in}}

% Slide content indentation
\newcommand{\contentstart}{%
\leftskip 0.6in \rightskip 0.6in}
\newcommand{\contentend}{%
\leftskip 0in \rightskip 0in}

\begin{document}
\begin{huge}

%%%%%%%%%%%%%%%%%%%%%%%%%%%%%%%%%%%%%%%%
% Slide 1 - Title
\slideheader{ECE5110 Numerical Modeling}{Unit 03}

\vspace{1.2in}
\begin{center}
{\Huge \textbf{Three-Point Finite-Difference Differentiation}}\\[0.3in]
{\Large Analytic Validation \& Free-Fall Gravity Estimation}
\end{center}

%%%%%%%%%%%%%%%%%%%%%%%%%%%%%%%%%%%%%%%%
\newpage
\slideheader{Motivation \& Goals}{Unit 03}
\contentstart
\begin{itemize}[leftmargin=*,itemsep=0.4em]
    \item Numerical differentiation is essential when analytic derivatives are unavailable.
    \item Three-point finite-difference stencils achieve $O(h^{2})$ accuracy with only three samples.
    \item This unit validates Python implementations on analytic benchmarks, then estimates $g$ from discrete free-fall data.
\end{itemize}
\contentend

%%%%%%%%%%%%%%%%%%%%%%%%%%%%%%%%%%%%%%%%
\newpage
\slideheader{Three-Point Differentiation Formulas}{Unit 03}
\contentstart
\begin{itemize}[leftmargin=*,itemsep=0.6em]
    \item \textbf{Central (symmetric):}
    \[ f'(x) \approx \frac{f(x+h)-f(x-h)}{2h} \]

    \item \textbf{Forward (one-sided):}
    \[ f'(x) \approx \frac{-3f(x)+4f(x+h)-f(x+2h)}{2h} \]

    \item \textbf{Backward (one-sided):}
    \[ f'(x) \approx \frac{f(x-2h)-4f(x-h)+3f(x)}{2h} \]
\end{itemize}
\contentend

%%%%%%%%%%%%%%%%%%%%%%%%%%%%%%%%%%%%%%%%
\newpage
\slideheader{Analytic Validation}{Unit 03}
\contentstart
\begin{itemize}[leftmargin=*,itemsep=0.4em]
    \item Three test functions with known derivatives were evaluated at fixed points.
    \item All 9 method–function combinations met the tolerances: central $1\times10^{-8}$, one-sided $1\times10^{-7}$.
    \item Example results (absolute error):
\end{itemize}

\vspace{0.2in}
\begin{tabular}{l c c}
\hline
Function & Method & Abs. Error \\
\hline
$\sin(x)$ @ $\pi/4$ & central & $1.4\times10^{-11}$ \\
$e^{x}$ @ 0.3 & forward & $3.4\times10^{-11}$ \\
$p(x)=x^{3}-2x^{2}+x-5$ @ 1.2 & backward & $1.4\times10^{-10}$ \\
\hline
\end{tabular}
\contentend

%%%%%%%%%%%%%%%%%%%%%%%%%%%%%%%%%%%%%%%%
\newpage
\slideheader{Free-Fall Gravity Estimation}{Unit 03}
\contentstart
\begin{itemize}[leftmargin=*,itemsep=0.4em]
    \item Seven position–time samples were collected during a drop test.
    \item Fit a quadratic polynomial $\hat y(t)=a_2 t^2 + a_1 t + a_0$ using \texttt{numpy.polyfit}.
    \item Apply central three-point differentiation twice (velocity then acceleration) at the mid-point $t_0$.
\end{itemize}

\vspace{0.3in}
\begin{center}
\includegraphics[width=0.78\linewidth]{../results/plots/freefall_gravity_interpolation.png}
\captionof{figure}{Quadratic interpolant fitted to measured free-fall samples (markers).}
\end{center}
\contentend

%%%%%%%%%%%%%%%%%%%%%%%%%%%%%%%%%%%%%%%%
\newpage
\slideheader{Gravity Estimate \& Results}{Unit 03}
\contentstart
\begin{itemize}[leftmargin=*,itemsep=0.4em]
    \item Estimated acceleration: $\hat a = -9.6904\ \mathrm{m/s^{2}}$ (downward sign).
    \item Magnitude $|\hat g| = 9.6904\ \mathrm{m/s^{2}}$.
    \item Comparison to standard $|g|=9.81\ \mathrm{m/s^{2}}$ yields error $0.1196\ \mathrm{m/s^{2}}$ ($1.2\%$).
    \item Result passes tolerance $0.15\ \mathrm{m/s^{2}}$.
\end{itemize}
\contentend

%%%%%%%%%%%%%%%%%%%%%%%%%%%%%%%%%%%%%%%%
\newpage
\slideheader{Key Takeaways}{Unit 03}
\contentstart
\begin{itemize}[leftmargin=*,itemsep=0.4em]
    \item Three-point stencils offer a simple, second-order-accurate way to approximate derivatives from sampled data.
    \item Central differences consistently produce smaller errors than one-sided formulas.
    \item Combining polynomial interpolation with numerical differentiation allows recovery of physical constants (e.g., $g$) from noisy, discrete measurements.
\end{itemize}
\contentend

%%%%%%%%%%%%%%%%%%%%%%%%%%%%%%%%%%%%%%%%
\newpage
\slideheader{Questions}{Unit 03}
\vspace{2.5in}
\begin{center}
{\Huge \textbf{Questions?}}
\end{center}

\end{huge}
\end{document}
